\documentclass[11pt,a4paper]{article}
\usepackage[T1]{fontenc}
\usepackage[utf8]{inputenc}
\usepackage{lmodern,amsmath,amssymb,amsthm,mathtools}
\usepackage[margin=24mm]{geometry}
\usepackage{booktabs,longtable,array,enumitem,microtype,xcolor,listings,hyperref}
\hypersetup{colorlinks=true,linkcolor=blue!45!black,urlcolor=blue!45!black,citecolor=blue!45!black,
 pdftitle={Smaller Diophantine descriptions of ZFC consistency},pdfauthor={Research contribution prepared for ProveIt}}
\setlength{\emergencystretch}{3em}
\setlength{\parskip}{3pt}
\setlist{itemsep=2pt,topsep=4pt}
\newtheorem{theorem}{Theorem}[section]
\newtheorem{lemma}[theorem]{Lemma}
\newtheorem{proposition}[theorem]{Proposition}
\newtheorem{corollary}[theorem]{Corollary}
\theoremstyle{remark}\newtheorem{remark}[theorem]{Remark}
\newcommand{\N}{\mathbb N}
\newcommand{\Z}{\mathbb Z}
\newcommand{\Pos}{\mathbb Z_{>0}}
\newcommand{\Con}{\operatorname{Con}}
\newcommand{\pc}{\operatorname{pc}}
\newcommand{\bits}{\operatorname{bits}}
\newcommand{\src}[1]{\texttt{\detokenize{#1}}}
\newcommand{\pin}{\texttt{3b201cdea32d2c816f5f3d20734051f25fd3ef14}}
\newcommand{\Parent}{\mathcal S_{\mathrm{norm}}}
\newcommand{\rgt}{\operatorname{right}}
\newcommand{\dn}{\operatorname{next}}
\lstset{basicstyle=\ttfamily\footnotesize,breaklines=true,columns=fullflexible,keepspaces=true,frame=single}
\title{\textbf{Smaller Diophantine descriptions of ZFC consistency}\\[5pt]
\large Thirteen positive unknowns, a smaller quartic,\\
and a four-site computation compiler}
\author{Research contribution prepared for the ProveIt repository}
\date{22 October 2026\\[3pt]\small Repository snapshot \pin}
\begin{document}
\maketitle
\begin{abstract}
We combine and extend the eight fixed-input constructions currently collected
in ProveIt's \emph{Con(ZFC) in Diophantine form} report. The main new equation
has thirteen strictly positive integer unknowns, total degree $571$, and an
explicit $84$-operation arithmetic circuit. Its unsolvability is equivalent,
provably in ZFC, to $\Con(\mathrm{ZFC})$. A four-equation presentation uses
$78$ operations. The key algebraic step replaces three square-root witnesses
by a single weighted sum, with a proof that the eliminated roots remain
integers and that the resultant has no degenerate denominator branch.

Other combinations improve different measures: a fifteen-variable,
six-equation system has degree $28$ and $61$ operations; its single
product has degree $99$ and $69$ operations. A direct lift gives $23$
quadratic equations in $32$ positive unknowns and hence one quartic. An
explicit domain refinement gives a quartic in $62$ unrestricted integer
unknowns. Finally, replacing nine-site windows by four-site windows reduces
the local alphabet to $71\,356\,645$ and the bit length of the distinguished
coefficient $B-1$ from $5^{67}$ to $5^{28}$, with the same four coefficient ports.

Every coefficient is fixed by a terminating finite recipe and the literal
K215 table. The enormous coefficients are not expanded. The new results
are ordinary mathematical proofs, accompanied by explicit circuits and
finite algebra and compiler checks; they are not newly checked Lean or
Rocq theorems. The constructions improve the reviewed collection in
specified measures, without claiming a global minimum, a new universal
operation bound, or a proof of ZFC consistency.
\end{abstract}
\clearpage
\tableofcontents
\clearpage

\section{Results and comparison}
\label{sec:results}

Write $\mathrm{Ref}_{\mathrm{ZFC}}$ for the arithmetic assertion that there
exists a finite derivation of a contradiction from ZFC. Throughout,
$\Con(\mathrm{ZFC})$ means $\neg\mathrm{Ref}_{\mathrm{ZFC}}$ for an ordinary
effective presentation. The machine calibration and its internalization
are discussed in Section~\ref{sec:logical}.

\begin{theorem}[Main result]\label{thm:main}
Let $(\beta,K_0,M_C,M_F)$ be the four fixed integers of
Section~\ref{sec:coefficients}, obtained from the literal blank-input K215
machine. For the polynomial $P_{13}$ defined in
Section~\ref{sec:thirteen},
\[
 \mathrm{ZFC}\vdash
 \bigl(\exists\mathbf x\in\Pos^{13}\ P_{13}(\mathbf x)=0\bigr)
 \ \longleftrightarrow\ \mathrm{Ref}_{\mathrm{ZFC}}.
\]
Consequently,
\[
 \boxed{\quad\mathrm{ZFC}\vdash
 \bigl(\forall\mathbf x\in\Pos^{13}\ P_{13}(\mathbf x)\ne0\bigr)
 \ \longleftrightarrow\ \Con(\mathrm{ZFC}).\quad}
\]
The same assertion holds for each alternative explicitly defined below,
in its stated domain. In particular there is a specified quartic in
$32$ positive unknowns and a specified quartic in $62$ unrestricted
integer unknowns with the corresponding equivalence.
\end{theorem}

The eight packages are numbered 01--08 in the repository's reconciliation
note~\cite{catalogue}. The best existing arithmetic baseline is the
independently obtained $15$-variable, $67$-operation polynomial of 06
and 08. Its degree is $155$. The four coefficients differ between those
packages; the same circuit shape does not make their numerals
interchangeable. We use 06's raw-K215 semantics as the base and prove a
new coefficient recipe for it.

\begin{center}\small
\begin{tabular}{lrrrr}
\toprule
Construction & Positive variables & Equations & Degree & Operations\\
\midrule
Existing 06/08 product &15&1&155&67\\
Existing 06 basic system &15&6&60&61\\
Existing 02 ordinary system &16&7&28&63\\
Existing 08 quadratic lift &34&25&2&---\\
\midrule
New ordinary system &15&6&28&61\\
New ordinary product $P_{15}$ &15&1&99&69\\
New retained-auxiliary system &16&7&24&60\\
New amalgamated system &14&5&72&67\\
New amalgamated product $P_{14}$ &14&1&121&76\\
New weighted system &13&4&544&78\\
New weighted product $P_{13}$ &13&1&571&84\\
New alternative-weight product &13&1&515&85\\
New quadratic lift &32&23&2&---\\
New quartic $Q_{32}$ &32&1&4&---\\
\bottomrule
\end{tabular}
\end{center}
For a system, ``degree'' means the maximum residual degree. Every
displayed degree is exact. Arithmetic counts charge every binary
$+$, $-$, or $\times$, including multiplication by a numeral. Reuse of
an already computed register and the final equality comparison are free.
The sum-of-squares alternatives are $78$ operations/degree $56$ for the
ordinary fifteen-variable system, $80$/degree $48$ for the retained
sixteen-variable system, $81$/degree $144$ for the fourteen-variable
system, and $89$/degree $1088$ for the thirteen-variable system.

There is a genuine tradeoff. The new thirteen-variable equation uses fewer
unknowns than all eight explicit packages, but does not beat the
$67$-operation evaluator. The new ordinary system lowers degree from
$60$ to $28$ at the same variable and operation counts. The quartic
reduces the positive-variable count from $34$ to $32$. In the unrestricted
integer domain our $62$-variable quartic improves the $83$-variable
quartic reported by 04; comparisons across different domains would be
misleading. Classical nine-variable existence results remain relevant
and prevent any claim that thirteen is an absolute minimum; the
repository's \src{Jones1982.nine_unknowns} is an existence theorem,
not an emitted polynomial with a stated degree~\cite{nine}.

\section{Common polynomials and the inherited interface}
\label{sec:common}

All free variables in Sections~\ref{sec:common}--\ref{sec:fourteen} are
strictly positive integers. A capitalized abbreviation does not introduce
an additional existential variable. Define
\begin{align}
 q&=\beta J+1,& C_0&=q-F-\alpha,&X&=wq,&Y&=sq^3,\notag\\
 E&=XY,& k&=\eta+\zeta,&c&=kY+\eta,&a&=E+Y,\notag\\
 H&=4a+3,&\Delta&=a^2+H=(a+2)^2-1,&
 D_0&=X+ac+\gamma H. \label{eq:common}
\end{align}
Put
\begin{align}
 G&=(q-1)(q-F)+\alpha,&
 R&=G(q^2-1)+(M_C+qM_F)J,\notag\\
 L_0&=E(kY),& A&=1+L_0(L_0+k), \label{eq:pack}\\
 N_1&=D_0^2-\Delta c^2,&
 N_k&=k-hE-R,\notag\\
 N_t&=(K_0+w)C_0+(q-F)-z(q-1). \label{eq:units}
\end{align}
Here $z$ is the transport quotient, while $T$ below is a different
auxiliary quotient. The coefficient $M_F$ is the \emph{shifted} source
mask, not the native one-cell high mask.

For comparison, package 06's normalized system uses the fifteen
unknowns
\[
 J,F,\alpha,z,f,h,i,T,s,w,\tau,\eta,\zeta,y,\gamma
\]
and, with
\[
 Q=(\Delta i c^2)^2,\qquad V=c(Tf-1)-Rf^2,
\]
has the six equations
\begin{equation}
 \tau^2=A,\quad N_1=N_k=N_t=1,\quad
 f^2=1+\Delta i^2c^4,\quad
 QV^2-(Q-1)y^2=1. \label{eq:normalized}
\end{equation}
This system is denoted $\Parent$. Package 06 proves that its positive
solvability is equivalent to the blank halting of K215~\cite{blank}.
Sections~\ref{sec:coefficients}--\ref{sec:compilerproof} prove that the
new four-site coefficients satisfy the same compiler interface, so the
same arithmetic theorem applies to them.

The following bounds will be used before any computation is decoded.
Suppose $N_t\in\{1,-1\}$. Since
\[
 N_t=(K_0+w)C_0+1-F-(z-1)(q-1),
\]
one has $C_0\ge0$; $C_0\le-1$ would make $N_t\le-2$.
One cannot yet assert $C_0>0$: the boundary
$C_0=0,F=2,z=1$ permits $N_t=-1$. With the authentic mask bounds,
\begin{gather}
 q\ge B\ge16,\qquad q\le G\le q^2-q,\notag\\
 0<(M_C+qM_F)J<(q-1)(1+2q),\notag\\
 q(q^2-1)<R<q^4-q^3+q^2-1<q^4-2,\label{eq:weakbounds}\\
 R>3q+1,\qquad E\ge q^4>R+2,\qquad a>R+2.\notag
\end{gather}
These follow directly from $q-F=C_0+\alpha>0$ and
$0<M_C<B-1$, $0<M_F<2(B-1)$, $J=(q-1)/(B-1)$.
If $N_k=\varepsilon\in\{1,-1\}$, then
$k=R+\varepsilon+hE>R$. The literal split of $c$ always gives
\begin{equation}
 kY<c<k(Y+1). \label{eq:ratio}
\end{equation}

\section{The thirteen-variable equation}
\label{sec:thirteen}

Supply exactly
\[
 \boxed{\quad J,F,\alpha,z,h,i,T,s,w,\eta,\zeta,\gamma,t\in\Pos.\quad}
\]
The three coordinates $f,\tau,y$ of \eqref{eq:normalized} are absent.
Using \eqref{eq:common}--\eqref{eq:units}, define
\begin{align}
 Q&=(\Delta i c^2)^2,&D&=1+\Delta i^2c^4,&v&=cT,&z_0&=c+RD,\notag\\
 d_0&=Q-1,&M&=D+v^2,&P&=QDv^2,&Z&=Qz_0^2,\notag\\
 e&=P+Z-1,&L&=d_0t^2,&N&=d_0M^2A,&C&=L+N-e,\notag\\
 a_0&=PZ,&b_0&=LN,&U&=C^2+4(a_0-b_0),\notag\\
 \mathcal R&=U^2-16C^2a_0.\label{eq:13defs}
\end{align}
The requested single equation can be taken to be
\begin{equation}
 \boxed{\quad P_{13}=N_1N_kN_t(1+\mathcal R)-1=0.\quad}
 \label{eq:P13}
\end{equation}
An equivalent existence condition is the four-equation system
\begin{equation}
 \boxed{\quad N_1=N_k=N_t=1,\qquad U^2=16C^2a_0.\quad}
 \label{eq:S13}
\end{equation}
All symbols in these displays are ordinary polynomial abbreviations.
There is no variable exponentiation, division, bit operation, primality
test, or implicit program parameter in the evaluated polynomial.
The proof uses rational expressions only to recover eliminated integer
witnesses.

\subsection{The weighted root lemma}
\begin{lemma}[Three roots represented by one integer]\label{lem:merge}
Let $A,D,Q,v,z_0,M$ be positive integers with $A,D,Q>1$, $v,z_0>0$,
and put $d_0=Q-1$, $e=Q(Dv^2+z_0^2)-1$.
Assume $d_0M^2A>e$. For $t>0$, set
\[
 C=d_0(t^2+M^2A)-e,\quad B_0=2Qvz_0,\quad
 U=C^2+DB_0^2-4d_0^2t^2M^2A.
\]
Then
\[
 U^2-D(2CB_0)^2=0
\]
implies the existence of nonzero integers $f,\tau,y$ satisfying
\[
 f^2=D,\quad\tau^2=A,\quad
 Q(vf-z_0)^2-(Q-1)y^2=1,\quad t=M\tau+y.
\]
Conversely any such roots with $t=M\tau+y>0$ give the resultant equation.
\end{lemma}
\begin{proof}
The hypothesis makes $C>0$, so $2CB_0\ne0$. Define
$f=-U/(2CB_0)$. The resultant gives $f^2=D$. A rational number whose
square is an integer is itself an integer: write it in lowest terms
and compare the denominator in the squared equation. Thus $f$ is an
integer and is nonzero. Expanding $U+2CB_0f=0$ gives
\[
 (C+B_0f)^2=4d_0^2t^2M^2A.
\]
Hence $\tau=(C+B_0f)/(2d_0tM)$ is a rational number with square $A$,
so it is a nonzero integer. Set $y=t-M\tau$. Then
\[
 d_0y^2=e-B_0f=Q(vf-z_0)^2-1.
\]
The value $y=0$ would imply that a multiple of $Q>1$ equals $1$.
The reverse implication follows by the same identities. No algebraic
number field or unproved rational-to-integer step is needed.
\end{proof}

\subsection{Why the chosen weight works}
At a zero of \eqref{eq:P13}, each of its four factors is an integer
unit. Since $\Delta$ is $0$ or $3$ modulo $4$, $N_1$ cannot equal
$-1$. Also $\mathcal R\equiv U^2\pmod4$, so $1+\mathcal R$ cannot
equal $-1$. Therefore $N_1=1$, $\mathcal R=0$, and
$N_k=N_t=\varepsilon\in\{1,-1\}$. The weak bounds above are
available at precisely this stage.

Write $U_0=XY^2$. Then
\[
 A=1+k^2U_0(U_0+1)>k^2Y^4>8k^2((Y+1)^2+1)>8(c^2+R^2).
\]
The middle inequality holds for $Y\ge4$, while here $Y\ge4096$.
For $M=D+v^2$,
\[
 M^2\ge4Dv^2,\qquad M^2\ge D^2,\qquad
 z_0^2=(c+RD)^2\le2D^2(c^2+R^2).
\]
Splitting $AM^2$ into two halves proves
$AM^2>2(Dv^2+z_0^2)$. Since $Q\ge2$ and $d_0\ge Q/2$,
\begin{equation}
 d_0AM^2>Q(Dv^2+z_0^2)=e+1.\label{eq:weight}
\end{equation}
Thus the lemma applies without already knowing the eliminated roots.
The weight is substantive: an unweighted elimination can have a
vanishing coefficient of $f$, so its resultant alone need not force
$D$ to be a square.

The shorter evaluator \eqref{eq:13defs} is the lemma's resultant because
$DB_0^2=4PZ$ and $d_0^2t^2M^2A=LN$.

\subsection{Soundness of the signed reconstructed roots}
The lemma restores \eqref{eq:normalized}, with $N_k,N_t$ still units
of equal sign and with $f,\tau,y$ provisionally signed. Replace
$\tau,y$ by their nonzero absolute values. If $f<0$, replace both
$f,T$ by their negatives; every relevant expression uses $f$ only
through $f^2$ and $Tf$. This makes $f>0$ while allowing $T$ to be
an arbitrary signed integer. It is essential to justify this domain
extension instead of applying a positive-$T$ theorem silently.

For an integer $A_0\ge2$, define the Pell sequences by
\[
 \chi_{A_0}(r)+\psi_{A_0}(r)\sqrt{A_0^2-1}
 =(A_0+\sqrt{A_0^2-1})^r.
\]
Use $A_0=a+2$ and $P_0=2XY^2+1$. The first and main norms imply
\[
 \tau=\chi_{P_0}(n),\quad k=2\psi_{P_0}(n),\quad
 D_0=\chi_{A_0}(p),\quad c=\psi_{A_0}(p).
\]
The index congruence gives $2n\equiv R+\varepsilon\pmod E$.
Since $0<R+\varepsilon<E$, it follows that
$n\ge(R-1)/2\ge24$. Since $P_0>A_0$ and $c>k$, Pell growth gives
$p>n$ and $c>2p$. Also $k>R$ implies $c>2R$.

The normalized strong equation gives
$f=\chi_{A_0}(m)$ and $\psi_{A_0}(m)=ic^2$. Strong divisibility of
$\psi$ implies $p\mid m$. Writing $m=pj$ and expanding
$(D_0+c\sqrt\Delta)^j$ modulo $c^2$ gives
\[
 \psi_{A_0}(pj)/c\equiv jD_0^{j-1}\pmod c.
\]
Because $\gcd(D_0,c)=1$, one obtains $c\mid j$; hence
$pc\mid m$ and $m>2p$.

Put $S=\Delta ic^2$ and retain the sign of $V$. The auxiliary norm
classifies $|SV|=\chi_S(\ell)$, $y=\psi_S(\ell)$ with odd
$\ell=2r+1$. If $\chi_X(2r+1)=X C_r(X^2)$, the recurrence gives
\[
 C_r(0)=(-1)^r\ell,\qquad C_r(-\Delta)=(-1)^r\psi_{A_0}(\ell).
\]
The literal congruences $V\equiv-c\pmod f$ and
$V\equiv-R\pmod c$, together with $S^2\equiv-\Delta\pmod f$,
give $R\equiv\pm p\pmod c$ by the signed Pell step-down argument.
For completeness, squaring the first congruence yields
\[
 \chi_{A_0}(2\ell)\equiv\chi_{A_0}(2p)\pmod{\chi_{A_0}(m)}.
\]
Reduction of the index modulo $2m$, reflection, and
$2\chi_{A_0}(m-1)<\chi_{A_0}(m)$ force
$\ell\equiv\pm p\pmod m$. The second congruence then gives the
claim modulo $c$. Since $0<R,p<c/2$, this forces $p=R$.
This argument uses $|V|$, and never assumes $T,V>0$
\cite{signed,scout}.

Now write $2n=R+\varepsilon+jE$ with $j\ge0$.
The bound $E>R+2$ excludes $j\ge1$, because $n<p=R$.
If $\varepsilon=-1$, then $p=2n+1$. Pell duplication, with
$\chi_{A_0}(2)>P_0$, would give
$c=\psi_{A_0}(2n+1)>A_0 k>k(Y+1)$, contrary to
\eqref{eq:ratio}. Thus $\varepsilon=1$, and $C_0>0$ follows.

The remaining exterior argument uses no sign of $V$ or $T$.
The lower Pell-ratio estimate gives $2^R<a$; the main-root
congruence modulo $H=4a+3$ gives $X\equiv2^R\pmod H$.
Since $0<X,2^R<a<H$, one has $X=2^R$. Hence $q\mid X$ makes
$q$ a power of two; the repunit relation makes $q=B^n$ for some
$n>0$. The bound $R>3q$ restores $q^3\mid X$.
The half-binomial extraction gives
\[
 \nu_2(Y)=\pc(R)-2,\qquad \pc(R)\ge3dn+2.
\]
The shifted-mask population identity supplies the reverse bound and
therefore the exact low and high masks. Their native convolution,
anchor alignment, and overlap constraints decode a blank halting
tableau. These exterior arguments are stated with their sign-independent
hypotheses in the inherited notes~\cite{scout,half,blank}; the new
tableau and coefficient interface is proved below. No numerical-input
Pell bridge is needed.

Conversely, a blank halting run supplies a positive solution of
\eqref{eq:normalized}. Set $t=M\tau+y>0$ in the lemma; then
\eqref{eq:S13} and \eqref{eq:P13} hold. This proves both directions.

\subsection{Degree and operation certificates}
\begin{center}\small
\begin{tabular}{lr@{\qquad}lr}
\toprule
Expression & Degree & Expression & Degree\\
\midrule
$c,a,\Delta,R$ &$5,6,12,4$&$A,Q,D,v,z_0$&$22,46,34,6,38$\\
$P,Z,M$&$92,122,34$&$L,N,C$&$48,136,136$\\
$a_0,b_0,U$&$214,184,272$&$\mathcal R$&$544$\\
$N_1,N_k,N_t$&$18,7,2$&$P_{13}$&$571$\\
\bottomrule
\end{tabular}
\end{center}
The highest homogeneous part of $C$ comes uniquely from $N$;
that of $U$ comes from $C^2$; and that of $\mathcal R$ comes from
$U^2$. All are nonzero for every $\beta>0$. For $N_1$, canceling
the apparent top degree in $D_0^2-\Delta c^2$ leaves the nonzero
degree-$18$ leading term $8\gamma a_{\rm top}^2c_{\rm top}$.
The product of nonzero leading forms in an integer polynomial ring
is nonzero. Hence $\deg P_{13}=18+7+2+544=571$ uniformly on the
actual coefficient slice.

The archive emits all $84=49M+35A$ live gates. Removing its final
unit-product gates and comparing $U^2$ with $16C^2a_0$ gives
$78=46M+32A$ operations for \eqref{eq:S13}. Adding four squared
residuals costs $11$ more operations and gives degree $1088$.
Modular univariate diagnostics supplement these symbolic leading-form
proofs; they are not the justification of uniform exact degree.

\subsection{One extra operation lowers the degree to 515}
There is a useful second weight. Put $I=ic^2$ and replace $M$ in
\eqref{eq:13defs} by
\[
 M=kI^2+T^2.
\]
All other polynomial definitions remain the same. For efficient evaluation,
compute $D=1+\Delta I^2$, $Q=\Delta(D-1)$, and $v^2=c^2T^2$.
These are identities with the earlier definitions, not new equations.
The resulting product uses $85=50M+35A$ operations and has degree
$515$; its four-equation system uses $79$ operations with maximum
degree $488$. The sum of squared residuals uses $90$ operations
and has degree $976$.

Here is the additional domination check. Write $A_0=a+2<2E$.
The weak bounds give $k>E$, $c<2kY$, $k>48$, and
\[
 L_0=kEY>A_0c/4,\qquad L_0>E^2Y>4A_0^2.
\]
Thus $A>A_0^2c^2/16$ and $A>16A_0^4$. Since $I>1$,
\[
 D<A_0^2I^2,\qquad M^2\ge4kI^2T^2,\qquad
 M\ge kI^2,\qquad z_0<2kA_0^2I^2.
\]
The first set of inequalities gives $AM^2>4Dv^2$; the second gives
$AM^2>4z_0^2$. Hence the same strict bound
$AM^2>2(Dv^2+z_0^2)$ proves $C>0$, so the whole restoration and
signed soundness argument applies unchanged. Now $\deg M=23$,
$\deg N=114$, while $e$ has degree $122$. The leading form of
$C$ is $-e_{\rm top}$; that of the resultant is $C_{\rm top}^4$.
Consequently its degree is $488$ and the product's degree is
$488+18+7+2=515$. The archive includes both weights; the first
has the shorter proof and cheaper evaluator.

\section{Lower degree with the ordinary strong equation}
\label{sec:ordinary}

Package 02's ordinary strong equation can be combined with 06's inputless
compiler. This changes the scaling of one witness, so the equivalence is
an equivalence of existence conditions, not a coordinatewise identity.
Retain the fifteen variables of \eqref{eq:normalized} and set
\[
 I=ic^2,\qquad K=\Delta(f^2-1),\qquad V=c(Tf-1)-Rf^2.
\]
The new system is
\begin{equation}
 \boxed{\quad
 \tau^2=A,\quad N_1=N_k=N_t=1,\quad
 I^2=K,\quad K(V^2-y^2)+y^2=1.\quad}
 \label{eq:ordinary}
\end{equation}
It has six equations, with exact residual degrees
$22,18,28,7,2,22$ in first/main/auxiliary/index/transport/strong order.
Its shared evaluator has $61=34M+27A$ operations. Define
\[
 N_0=\tau^2-L_0(L_0+k),\quad
 N_a=K(V^2-y^2)+y^2,\quad N_s=1+I^2-K.
\]
An equivalent single equation is
\begin{equation}
 \boxed{\quad P_{15}=N_0N_1N_aN_kN_tN_s-1=0.\quad}
 \label{eq:P15}
\end{equation}
It has $69=39M+30A$ operations and exact degree $99$.
The sum of the six squared residuals has $78$ operations and degree $56$.

\subsection{Why the relaxed equation is sufficient}
Completeness is immediate from the normalized system: replace its
$i$ by $\Delta i$. Then the new $I^2=\Delta(f^2-1)$ equals the
old auxiliary coefficient $(\Delta ic^2)^2$, preserving all other
coordinates and equations.

Soundness uses the following independently established interface
\cite{ordinary,relaxed}.
\begin{lemma}[Ordinary strong rank]\label{lem:ordinaryrank}
Let $A_0\ge2$, $\Delta=A_0^2-1$, $c=\psi_{A_0}(p)$, and
$c>A_0\Delta^2$. If positive integers $i,f$ satisfy
$(ic^2)^2=\Delta(f^2-1)$, there is $m>0$ such that
\[
 f=\chi_{A_0}(m),\quad ic^2=\Delta\psi_{A_0}(m),
 \qquad p\mid m,\quad c\mid m.
\]
In particular $c^2\mid f^2-1$.
\end{lemma}
\begin{proof}[Proof mechanism]
Factor $\Delta$ into its squarefree part times a square and use the
fundamental norm-one unit in the integer-coefficient quadratic ring.
If $A_0=\chi_F(e)$ and $S=\psi_F(e)$, then
$\Delta=(F^2-1)S^2$ and $\psi_F(ep)=Sc$.
For the auxiliary index $n$, strong divisibility gives
$h=\psi_F(\gcd(n,ep))\ge c/\Delta$. If $\gcd(n,ep)<ep$,
the proper-divisor bound and Pell doubling give
$h^2<A_0c/2$, contradicting $c>A_0\Delta^2$.
Thus $ep\mid n$, recovering an integral Pell solution at $A_0$.
Writing $m=pj$, the binomial expansion modulo $c$ yields
$c\mid\Delta j$. Finally
$\gcd(c,\Delta)=\gcd(p,\Delta)\mid p$, so $c\mid pj=m$.
The full elementary proof, including squarefree divisibility and the
fundamental-unit argument, is included in the archive~\cite{relaxed}.
\end{proof}

The first/main norms and the index congruence establish
$p>n\ge24$, $c>(a+2)\Delta^2$, $c>2p$, and $c>2R$
before this lemma is used. This ordering matters: the individual relaxed
norm can have extra rational Pell solutions without these hypotheses.
The lemma gives $m\ge c>2p$. It also gives $f>2c$ and
\[
 K-(Rf^2+c)=(\Delta-R)f^2-\Delta-c>0.
\]
The auxiliary norm excludes $V=0$. If $|V|>1$, it implies
$|V|\ge2K-1$ by comparing consecutive squares. But $T>0$ gives
$V>-Rf^2-c>-K$, excluding negative $V$ of magnitude greater than
one. The congruence $V\equiv-c\pmod f$, with $f>2c$, excludes
$V=\pm1$. Hence $V>1$.

The integers
\[
 o=cT-Rf,\qquad j=Tf-1-R(f^2-1)/c
\]
are positive because $of=V+c$ and $jc=V+R$. The quotient in the
definition of $j$ is integral by Lemma~\ref{lem:ordinaryrank}.
The two congruences $V=of-c=jc-R$ give $p=R$ by the same auxiliary
step-down, now with coefficient $I^2$. The exterior power, population,
and tableau decoding follows as before.

For \eqref{eq:P15}, first infer integer units. The consecutive-product
negative Pell descent excludes $N_0=-1$; modulo $4$ excludes
$N_1=-1$. Since $K=\Delta(f^2-1)$ is $0$ or $1$ modulo $4$,
$N_s=1+I^2-K$ cannot be $-1$. Thus $K=I^2$, and $N_a$ is also
a square modulo $4$, so $N_a=1$. The remaining equal signs on
$N_k,N_t$ are resolved with the weak $C_0\ge0$ bounds, ordinary rank,
and the same index-sign argument used above. This proves product
soundness without prematurely assuming $C_0>0$.

\subsection{Retaining two auxiliary witnesses}
Replace $T$ by positive $o,j$, put $V=of-c$, and append
$V=jc-R$ to \eqref{eq:ordinary}. This gives sixteen variables and
seven equations. The residual degrees are $22,18,24,7,2,22,6$;
the evaluator costs $60=33M+27A$ operations. Its soundness proof is
shorter: once $f>2c$, the inequality $V=of-c>c>1$ is immediate.
The sum of squares has degree $48$ and $80$ operations. These
retained witnesses provide the useful starting point for the quartic.

\section{A fourteen-variable intermediate option}
\label{sec:fourteen}

The ordinary system also permits a less costly root combination.
Remove $\tau,y$ and retain $f$, supplying one positive integer $t$.
Let $I,K,V,A$ have their ordinary meanings and put
\[
 d=K-1,\quad e=KV^2-1,\quad L=dt^2,\quad N=dA,\quad
 u=L+N-e,\quad \mathcal R_{14}=u^2-4LN.
\]
The five equations are
\[
 N_1=N_k=N_t=1,\qquad I^2=K,\qquad u^2=4LN.
\]
They use $67=38M+29A$ operations and have maximum degree $72$.
At a solution $K=I^2>1$, and
\[
 \tau=\frac{u}{2dt},\qquad \tau^2=A,\qquad y=t-\tau
\]
recover integral nonzero $\tau,y$ and $dy^2=e$.
Their absolute values give an ordinary-system solution. Conversely,
$t=\tau+y$ gives the five equations.

The product
\[
 \boxed{\quad P_{14}=N_1N_kN_t(1+I^2-K)(1+\mathcal R_{14})-1=0\quad}
\]
has $76=42M+34A$ operations and degree $121$.
Both new unit factors exclude $-1$ modulo $4$; therefore the same
reconstruction and ordinary sign proof apply. Its sum-of-squares
alternative has $81$ operations and degree $144$.
The archive additionally records the inputless specialization of 05's
$f$-elimination: fourteen-variable products at $73$ operations/degree
$342$ and $75$/degree $293$. These offer other arithmetic tradeoffs,
but are not the main result here.

\section{An explicit quartic in 32 positive unknowns}
\label{sec:quartic}

Retain the sixteen primary unknowns
\[
 J,F,\alpha,z,f,h,i,o,j,s,w,\tau,\eta,\zeta,y,\gamma
\]
and adjoin sixteen positive register unknowns
\[
 N,E,P,c,a,X,D_0,C_2,G,L,F_2,V,g,I,K,\Delta.
\]
Only $q=\beta J+1$, $k=\eta+\zeta$, $H=4a+3$, and
$C_0=q-F-\alpha$ remain abbreviations. They are affine in the
unknowns. Here the symbol $L$ means $q^2-1$, and $P=L_0+k$;
these register names are local to this section.

The following are the complete $23$ quadratic equations:
\begin{align}
 L&=q^2-1,&N&=q(L+1),\notag\\
 X&=wq,&a&=E+sN,\notag\\
 E&=X(a-E),&c&=k(a-E)+\eta,\notag\\
 P&=E(c-\eta)+k,&\tau^2&=P(P-k)+1,\notag\\
 D_0&=X+ac+\gamma H,&\Delta&=a^2+H,\notag\\
 C_2&=c^2,&D_0^2&=\Delta C_2+1,\notag\\
 k&=jc-V+1+hE,&G&=(q-1)(q-F)+\alpha,\notag\\
 jc-V&=LG+(M_C+qM_F)J,&&\notag\\
 (K_0+w)C_0+(q-F)-z(q-1)&=1,&&\notag\\
 F_2&=f^2,&V&=of-c,\notag\\
 g&=y^2-V^2,&I&=iC_2,\notag\\
 K&=I^2,&K&=\Delta(F_2-1),\notag\\
 y^2-Kg&=1.&&\label{eq:quad23}
\end{align}
Let $r_1,\ldots,r_{23}$ be their left sides minus right sides. Then
\begin{equation}
 \boxed{\qquad Q_{32}=\sum_{j=1}^{23}r_j^2=0.\qquad}
 \label{eq:quartic32}
\end{equation}
This is a single polynomial of exact total degree $4$. Only the first
Pell residual contains $\tau$, with coefficient $1$ on $\tau^2$,
so the coefficient of $\tau^4$ in $Q_{32}$ is $1$.

To project a solution, set $Y=a-E=sq^3>0$ and $R=jc-V$.
The packing equation forces the original expression for $R$; the
transport equation forces $C_0>0$, hence $R>0$.
Every equation of the retained-auxiliary system is recovered.
Conversely its positive solutions extend to all sixteen registers.
For the less immediate sign, the auxiliary norm gives
\[
 (K-1)(y^2-V^2)=V^2-1>0,
\]
so $g>0$. All other register signs follow from their displayed
definitions and the earlier budget bounds. Thus positive solvability
is equivalent in both directions.

The two-register saving over 08 has a specific explanation.
Ordinary scaling removes the separate register for $\Delta c^2$ used
by the normalized strong multiplier. Replacing $T$ by $o,j$ costs
one primary variable but removes the register $Tf-1$. Finally,
$R=jc-V$ removes the packed-index register and its defining equation
without increasing any residual's degree. The final count is $32$
unknowns and $23$ quadratics instead of $34$ and $25$.

\section{Domains: natural numbers and unrestricted integers}
\label{sec:domains}

\subsection{Natural-number forms}
To obtain a polynomial over $\N=\{0,1,2,\ldots\}$, replace every
positive coordinate $x$ by $x+1$. This preserves the number of
unknowns and the exact degree. It changes the displayed evaluation
cost if the extra additions are charged; no positive-domain count
is silently reused for the shifted evaluator.

\subsection{A quartic in 62 unrestricted integer unknowns}
In \eqref{eq:quad23}, positivity need be explicitly imposed only on
the following ten primary coordinates:
\begin{equation}
 \mathcal P=\{J,F,\alpha,z,i,s,w,\eta,\zeta,\gamma\}.
 \label{eq:tenpositive}
\end{equation}
The six coordinates $f,h,o,j,\tau,y$ and all sixteen registers
can be allowed to range over $\Z$. This relaxation requires a
soundness proof; it is not a rule that arbitrary auxiliary signs
may be discarded.

With \eqref{eq:tenpositive}, the quadratic definitions force
$q,X,Y,E,k,c,a,\Delta,D_0,I,K>0$ and $K>1$.
The transport equation forces $C_0>0$, giving $0<R<E-2$.
Since $k=R+1+hE>0$, an integer $h\le-1$ is impossible.
Thus $h\ge0$ and $k>R$, which are all that the rank bootstrap needs.
The first norm makes $\tau\ne0$, and ordinary strong scaling
makes $|f|>1$. Apply Lemma~\ref{lem:ordinaryrank} to $|f|$.
The auxiliary norm makes $y,V\ne0$ and gives the two congruences
\[
 V\equiv-c\pmod{|f|},\qquad V\equiv-R\pmod c.
\]
Classifying the absolute auxiliary Pell root gives $p=R$ with
exactly the signed step-down used in Section~\ref{sec:thirteen}.
Neither $o$ nor $j$ needs a sign in these congruences. The no-wrap
argument gives $2n=R+1$; in fact $k=2\psi_{P_0}(n)>2n=R+1$
then restores $h>0$.

For exterior decoding, use only the forward ratio, half-binomial
extraction, and valuation statements that depend on $p=R$,
$2n=R+1$, and positive $X,Y,c,k$. They do not require $V>0$ or a
prior assertion $R\equiv3\pmod4$. The resulting population threshold
forces the native masks and hence the same blank tableau. One must not
invoke the complete positive-domain kernel on this signed tuple
without this qualification. Positive completeness is unchanged.

Every positive integer has the division-free representation
\begin{equation}
 \operatorname{Pos}(u,v,w)=1+u^2+v^2+w(w+1),\qquad u,v,w\in\Z.
 \label{eq:pos}
\end{equation}
Its values are positive. For surjectivity, the three-square theorem
applied to $4n-3\equiv1\pmod4$ gives one odd square and two even
squares; write them as $(2w+1)^2,(2u)^2,(2v)^2$.
For each $x\in\mathcal P$, append the quadratic equation
\[
 x-\operatorname{Pos}(u_x,v_x,w_x)=0.
\]
There are $32+3\cdot10=62$ unrestricted integer unknowns and
$23+10=33$ quadratic equations. Define
\begin{equation}
 \boxed{\quad Q_{62,\Z}=\sum_{j=1}^{23}r_j^2+
 \sum_{x\in\mathcal P}
 \bigl(x-1-u_x^2-v_x^2-w_x(w_x+1)\bigr)^2=0.\quad}
 \label{eq:quartic62}
\end{equation}
This is an ordinary integer polynomial of exact degree $4$; the
$\tau^4$ coefficient is still $1$. The preceding signed-domain
argument proves
\[
 \mathrm{ZFC}\vdash
 \bigl(\neg\exists\mathbf z\in\Z^{62}\ Q_{62,\Z}(\mathbf z)=0\bigr)
 \longleftrightarrow\Con(\mathrm{ZFC}).
\]

\subsection{A 39-variable integer version of the main equation}
Substituting a separate copy of \eqref{eq:pos} for each of the
thirteen positive coordinates of $P_{13}$ gives a specified
polynomial in $39$ unrestricted integer variables. Its degree is
exactly $1142$, and its evaluator has $84+13\cdot7=175$
operations. Each substitution uses three multiplications and four
additions. Nonvanishing of the highest homogeneous part follows
by setting the second and third coordinates of each block to zero;
the remaining substitution $x_j\mapsto u_j^2$ is injective on
monomials. This offers fewer unrestricted unknowns than the
$41$-variable integer products of 06, at a larger degree.
\section{Fixed coefficients from four-site windows}
\label{sec:coefficients}

This section specifies the four integers used by every new construction.
Only coefficient construction uses finite enumerations, exponentiation,
and bit lengths. Once constructed, these are fixed integers, not extra
unknowns or additional operations in the displayed polynomial circuit.

\subsection{The exact machine and tile relation}

Use the literal \src{compiler/k215.tm} supplied with this archive. Its
SHA-256 digest is
\[
\texttt{fdbcef0825886035c3091c8537a69535bc96f9e265f0ae45863543d40ae70090}.
\]
The source has $215$ working states $0,\ldots,214$, binary symbols,
and halt state $215$. Complete the halt state by staying put and
preserving its scanned symbol. The first transition is
$(0,0)\mapsto(93,1,\mathrm{right})$.
This is the raw machine calibrated in the K215 endpoint, with no added
input or initialization machine~\cite{endpoint,blank}.

Let $a_t=439$. The tile alphabet has headless bits $0,1$, head-bearing
tiles $2(s+1)+b$ for $0\le s\le215$ and $b\in\{0,1\}$, and
\[
 L=434,\quad Q_{\rm init}=435,\quad R_{\rm init}=436,\quad
 V_{\rm bdry}=437,\quad H_{\rm bdry}=438.
\]
The three phase tiles have blank payloads, with an initial state-$0$
head only at $Q_{\rm init}$. $V_{\rm bdry}$ has the vertical boundary
flag; $H_{\rm bdry}$ has the horizontal boundary flag. All other
tiles have neither flag.

The original nine-site local relation inspects only its cross
$(N,W,C,E,S)$. Moreover, its only use of $N$ is the initialization
presence test: an interior center has a phase if and only if $N$
is a horizontal boundary. Translate that test one time step forward:
an interior $S$ has a phase if and only if $C$ is a horizontal
boundary. All other rules can be stated on the four-site tuple
$(W,C,E,S)$:
\begin{enumerate}
\item The vertical flag of $C$ equals that of $S$. Two horizontally
adjacent vertical boundaries and two temporally adjacent horizontal
boundaries are forbidden. Away from vertical boundaries the horizontal
flag propagates to the right.
\item If $S$ is interior, its phase is present exactly when $C$ is a
horizontal boundary. After this test, a boundary center has no
further center rule.
\item At a phase center, a left vertical boundary requires phase $L$,
and a right vertical boundary requires phase $R_{\rm init}$.
If the right neighbor is not a vertical boundary, the phase pair is
one of
$(L,L),(L,Q_{\rm init}),(Q_{\rm init},R_{\rm init}),
(R_{\rm init},R_{\rm init})$.
\item A head immediately before a horizontal boundary must be halted.
A moving head cannot escape through a neighboring vertical boundary.
\item When $S$ is not a horizontal boundary, its payload is the
deterministic successor of the three payloads at $W,C,E$.
Vertical boundary neighbors count as headless blanks; other neighbors
must have payloads. A center writes according to its transition;
incoming heads from either neighbor and a staying center head are
combined only if at most one head arrives.
\end{enumerate}
The executable predicate is \src{compiler/tshape_helical.py:local_valid}.
It implements exactly these clauses, using the literal completed table.

\begin{lemma}[Translation of one local test]\label{lem:translate}
A tile plane satisfies the four-site relation at every center if and
only if it satisfies the original nine-site relation at every center.
The statement holds on the full integer plane and on every periodic
plane, including presentations whose spatial and temporal shifts
are not independent.
\end{lemma}
\begin{proof}
The four corners of the original window are unused. All remaining
clauses are unchanged except for the single presence test. Temporal
translation is a bijection of either plane, so requiring that test
at every center is equivalent to requiring its translate at every
center. No primitivity or distinct-offset assumption enters.
\end{proof}

The distinguished Start tuple and the excluded tuple are
\[
 U_0=(434,435,436,1),\qquad U_1=(437,438,437,435).
\]
Both satisfy their own local predicate. Start is forced at every
initialized head, by its phase neighbors and the literal first
transition. The excluded tuple has no global extension: its left
vertical boundary propagates to the following row, where the
following initialized $Q_{\rm init}$ would have a vertical boundary
on its left, violating the phase endpoint rule.

\subsection{Exact local alphabet count}

There are $349$ symbols sending a head from the left, including
$Q_{\rm init}$, and $82$ sending one from the right. Among the $438$
non-horizontal-boundary symbols this leaves $89$ and $356$
nonincoming choices, respectively. The following census is exhaustive.
\begin{center}\small
\begin{tabular}{lrr}
\toprule
Center type & Multiplicity & Valid tuples per center\\
\midrule
Vertical boundary &1&192\,282\\
Horizontal boundary &1&2\,634\\
Headless bit &2&355\,508\\
Halted head &2&223\,966\\
Working head &430&162\,788\\
Initial $L$ &1&1\,754\\
Initial $Q_{\rm init}$ &1&437\\
Initial $R_{\rm init}$ &1&1\,750\\
\bottomrule
\end{tabular}
\end{center}
For example the headless count is
\[
 439\cdot438+89\cdot356+349\cdot356+82\cdot89.
\]
The first term ends the row at a horizontal boundary; the other terms
allow zero or one incoming head and then determine the next payload.
For a working head, its escape direction excludes one boundary
choice. The left-moving count is
\[
 (89-1)356+349\cdot356+82(89-1)=162\,788,
\]
and the right-moving count has the same value. The phase pair graph
and the endpoint rules give the last three rows. Thus the alphabet
of valid four-site tuples has exactly
\[
 \boxed{\ \kappa=71\,356\,645.\ }
\]
The archive independently checks the formulas against the actual
predicate for every left/right pair in each of thirteen exhaustive
center categories. Enumerate valid tuples with $U_0$ first,
$U_1$ second, and all remaining tuples in lexicographic order.
This defines their selector numbers without any arbitrary choice.

\subsection{Three overlaps suffice}

For selected tuples $U_i$, indexed cyclically by $\Z/N\Z$, impose
\begin{equation}
 U_i.E=U_{i-1}.C,\qquad
 U_i.C=U_{i-1}.W,\qquad
 U_i.S=U_{i-h}.C. \label{eq:overlaps}
\end{equation}
Putting $a_i=U_i.C$ recovers
\[
 U_i=(a_{i+1},a_i,a_{i-1},a_{i-h}).
\]
The pullback plane at spatial/time coordinates $(x,y)$ is
$a_{-x-hy}$. Every actual neighborhood is therefore valid. This works
for arbitrary $N,h$, with no assumption that $h\mid N$ or
$\gcd(N,h)=1$. Start rules out $h=0$ modulo $N$ because its center
and following symbols differ.

Lemma~\ref{lem:translate} now gives the same finite-strip soundness
as 06. Scanning horizontally from an initialized head finds vertical
boundaries on both sides; the phase chains cannot pass through that
head and close. The boundary tracks persist in time. In the finite
cyclic time orbit, an initializing horizontal boundary eventually
recurs; the first recurrence ends a genuine bounded blank computation,
whose last head must be halted. Conversely a finite halting run admits
the independently padded helical presentations of 06. Their Start
and absence of the excluded selector survive the projection
\cite{blank,helical}.

\subsection{Complete finite recipe for the four integers}

Use the spatial payload set and copy-target sets
\[
 \mathcal P=((1,0),(1,1),(1,2),(2,1)),\qquad
 \mathcal V=((1,1)),\qquad
 \mathcal H=((1,0),(1,1)).
\]
Here $\mathcal P$ is ordered as written. For a positive integer $x$,
$\bits(x)$ is its binary length. Put
\begin{align*}
 \ell&=\max(2,\bits(\kappa+1)),& A_b&=2^\ell,\\
 p_d&=\max\left(0,\left\lceil
            \frac{\kappa-2\ell-7a_t-4}{2}\right\rceil\right),&
 j_{\max}&=4a_t+4+p_d,\\
 \mu&=A_b-2+\sum_{j=1}^{j_{\max}}A_b^j,&
 m_0&=2\ell+11a_t+8+2p_d,\\
 n_{\rm pos}&=m_0+1,&
 n_{\rm dum}&=n_{\rm pos}-\kappa-4a_t-4,\\
 e_*&=\kappa+12a_t,&
 M_0&=\kappa+15a_t+n_{\rm dum},\\
 E_{\max}&=27M_0,&
 H_c&=E_{\max}+24M_0+3a_t+1,\\
 T_1&=H_c+2E_{\max}+a_t+1,&
 T_2&=T_1+4E_{\max}+2.
\end{align*}
Let $L_{\rm cell}$ be the least power of $5$ at least
$T_2+2E_{\max}+2$. For $(r,c)\in\mathcal P$ and $0\le u<a_t$, put
\begin{align*}
 e(r,c,u)&=\kappa+3a_t+(2-r)3a_t+(2-c)a_t+u,\\
 j(r,c,u)&=1+a_t\,\operatorname{index}_{\mathcal P}(r,c)+u,
\end{align*}
with indices starting at zero. Form the disjoint set $\mathcal E$
of native positions from the $\kappa$ selectors $0,\ldots,\kappa-1$,
the $4a_t$ retained payload positions, $n_{\rm dum}$ consecutive
dummy positions beginning at $e_*$, and anchors
$M_0,3M_0,9M_0,27M_0$.

At each $e\in\mathcal E$ give a clause coefficient $c_e$ as follows:
\begin{align*}
 c_i&=1+\sum_{(r,c)\in\mathcal P}A_b^{j(r,c,U_i[r,c])}
               +\sum_{t=1}^{4}A_b^{4a_t+t}
       &&(0\le i<\kappa),\\
 c_{e(r,c,u)}&=A_b^{j(r,c,u)},&
 c_{3^{t-1}M_0}&=A_b^{4a_t+t}\quad(1\le t\le4),\\
 c_e&=0&&\text{at a dummy position}.
\end{align*}
Here $U_i[r,c]$ is its corresponding $W,C,E,S$ component.
Define
\begin{align*}
 U_{\rm bd}&=4(n_{\rm pos}+2)\bigl(28(\kappa+4a_t+4)+6\bigr)+8,\\
 b_{\min}&=\max\bigl(\ell(j_{\max}+1)+3,\,
              \ell(4a_t+4)+\bits(U_{\rm bd})+1,\,4\bigr).
\end{align*}
Let $b$ be the least power of $5$ at least $b_{\min}$ and put
\[
 \rho=2^b,\qquad d=bL_{\rm cell},\qquad B=\rho^{L_{\rm cell}}=2^d.
\]
Finally set
\begin{align}
 D_C={}&\rho^{3a_t}+\rho^{H_c+a_t}+\rho^{8M_0}+\rho^{24M_0}
       +\sum_{e\in\mathcal E}c_e
          \bigl(\rho^{T_1-e}+\rho^{T_2-e}\bigr),\notag\\
 D_R={}&\rho^{H_c},\notag\\
 M_{F,0}={}&\mu(\rho^{T_1}+\rho^{T_2})
       +\sum_{\substack{(r,c)\in\mathcal V\\0\le u<a_t}}\rho^{e(r,c,u)}
       +\sum_{\substack{(r,c)\in\mathcal H\\0\le u<a_t}}
                      \rho^{H_c+e(r,c,u)}
       +\rho^{9M_0}+\rho^{27M_0},\notag\\
 M_C={}&B-1-\sum_{\substack{e\in\mathcal E\\e\ne1}}\rho^e-2\rho^{e_*},
       \qquad M_{F,\mathrm{native}}=M_{F,0}+4. \label{eq:coeffrecipe}
\end{align}
The promised four ports are
\begin{equation}
 \boxed{\quad
 \beta=B-1,\qquad K_0=D_C+BD_R,\qquad
 M_C,\qquad M_F=M_{F,\mathrm{native}}+B-1.\quad}
 \label{eq:ports}
\end{equation}
This is a closed specification of four particular natural numbers.
It uses the complete finite transition table and finite alphabet only;
it does not run K215 or consult whether ZFC is consistent.

For clarity, the geometry evaluates to
\begin{center}\small
\begin{tabular}{lr@{\qquad}lr}
\toprule
Quantity & Exact value & Quantity & Exact value\\
\midrule
$\ell$&27&$p_d$&35\,676\,757\\
$j_{\max}$&35\,678\,517&$m_0$&71\,358\,405\\
$n_{\rm pos}$&71\,358\,406&$n_{\rm dum}$&1\\
$M_0$&71\,363\,231&$E_{\max}$&1\,926\,807\,237\\
$H_c$&3\,639\,526\,099&$T_1$&7\,493\,141\,013\\
$T_2$&15\,200\,369\,963&$b_{\min}$&963\,319\,989\\
$b$&$5^{13}$&$L_{\rm cell}$&$5^{15}$\\
\bottomrule
\end{tabular}
\end{center}
In particular,
\begin{equation}
 \boxed{\ d=5^{28}=37\,252\,902\,984\,619\,140\,625.\ }
 \label{eq:newd}
\end{equation}
The exact inequalities
\[
 5^{12}<963\,319\,989\le5^{13},\qquad
 5^{14}<19\,053\,984\,439\le5^{15}
\]
verify the two least-power choices. The integers are intentionally
specified by recipe: expanding even $B-1$ would require $d$ bits.

\section{Why the smaller coefficient recipe is sound}
\label{sec:compilerproof}

Merely checking mask congruences is insufficient to replace the
compiler constants. We verify the actual supports, clause fields,
alignment, and completeness hypotheses of 06. The full parent
compiler proof and helical proof are included under \src{upstream};
the following identifies every changed component.

\subsection{Clause matching, copy fields, and mask balance}

The payloads retain their former spatial exponents. At a vertical
target $e(r,c,u)$, multiplication by $\rho^{3a_t}$ selects exactly
$e(r+1,c,u)$. At a horizontal target $H_c+e(r,c,u)$,
multiplication by $\rho^{H_c+a_t}$ selects exactly $e(r,c+1,u)$,
whereas $D_R$ selects the next window's $e(r,c,u)$. All three
source/target pairs belong to $\mathcal P$. Thus these are precisely
the three equalities in \eqref{eq:overlaps}. The selectors stop below
$\kappa$; payloads stop below $\kappa+12a_t$; dummies lie above the
payloads. None creates an extra source at a retained copy target.

Compact clause indices remove only unused fields. A selected tuple
supplies one selector, its four matching tile indicators, and the
four anchors. The same two repeated clause bands therefore enforce
exact selector/payload agreement. There are $3a_t$ copy parity tests,
and
\[
 \pc(\mu)=\ell+4a_t+3+p_d,\qquad
 \pc(M_{F,0})=2\pc(\mu)+3a_t+2=m_0.
\]
The native-position census $n_{\rm pos}=m_0+1$ gives exactly
\begin{equation}
\begin{gathered}
 \pc(M_C)+\pc(M_{F,\mathrm{native}})=d,\qquad
 0<M_C,M_{F,\mathrm{native}}<B-1,\\
 M_C\equiv2\pmod4,\qquad
 M_{F,\mathrm{native}}\equiv4\pmod8.
\end{gathered}\label{eq:maskproperties}
\end{equation}
These are equalities about the newly constructed masks, not
hypotheses checked only on example masks.

\subsection{Carry control and synchronization}

The coefficient mass estimate used in the parent proof becomes
\[
 \sum_{e\in\mathcal E}c_e
 \le14(\kappa+4a_t+4)A_b^{4a_t+4}.
\]
There are at most $\kappa$ selector terms, $4a_t$ payload terms,
and four anchor terms; bounding each by its positive geometric
sum gives the displayed estimate. The choice of $b_{\min}$
supplies the same quarter-radix carry margin as in 06.

Every native position is in $[0,E_{\max}]$. The inequalities
\[
 T_2-T_1>2E_{\max},\quad
 L_{\rm cell}-(T_2-T_1)>E_{\max},\quad
 T_2+E_{\max}<L_{\rm cell}
\]
separate the two clause bands and their rotations. Also $e_*+1<E_{\max}$.
The clean-band argument consequently synchronizes a transport field
to a whole-cell temporal shift before any copy equality is decoded.
The anchor pair $9M_0,27M_0$ still has the required unique difference
$18M_0$. Omitting corner and north payloads changes neither assertion.
A selected tuple has one tile indicator at every retained position;
an empty cell has none. The horizontal $E/C$ copy equality propagates
occupancy from the marked Start throughout the cyclic word.

The sparse checker evaluates these support and coefficient statements
for every native basis bit at every tested position in its finite
fixtures. The general argument above supplies their parameter-uniform
justification; small fixtures do not certify an enormous enumeration.

\subsection{Actual-index completeness}

The five-adic correction remains exactly
\[
 2D_C-D_R\equiv3\pmod5.
\]
Indeed $T_2-T_1=4E_{\max}+2$ makes the two clause-band sums cancel
modulo $5$. The four unpaired monomials, $a_t=439\equiv3\pmod4$,
and $\rho\equiv2\pmod5$ give the displayed residue.

Hence the parent's ignored-dummy argument still adjusts the
\emph{actual final} packed index to the desired temporal stride.
The one dummy position supplies an independently selectable bit in
each word cell. Its Boolean subset argument uses the same unit
coefficient modulo the five-power padding modulus. Canonical
helical tableaux have one Start and nonzero temporal stride; these
ensure the exceptional origin mask test. The dummies remain above
the low selector fields. The same canonical carry estimate gives
$2C_0+F<q/3$, so $\alpha=q-C_0-F>0$.
Only after this adjustment is the positive Pell converse invoked,
at that final index. This avoids assuming a converse at an index
that was later changed.

Thus both directions of the genuine native-compiler interface hold
with \eqref{eq:ports}. Composing it with the arithmetic proofs
establishes the new constructions. The intermediate five-site
cross compiler, also included in the archive, needs no translated
test and gives $\kappa=38\,256\,607\,897$ and $d=5^{36}$.
The original nine-site alphabet is exactly the cross alphabet
times $439^4$, because its four corners are unconstrained.
\begin{center}\small
\begin{tabular}{lrrr}
\toprule
Window & Local alphabet size & Payload groups & $\bits(B-1)$\\
\midrule
Original nine-site &1\,420\,903\,358\,357\,108\,792\,377&9&$5^{67}$\\
Five-site cross &38\,256\,607\,897&5&$5^{36}$\\
New four-site &71\,356\,645&4&$5^{28}$\\
\bottomrule
\end{tabular}
\end{center}
The final reduction in the distinguished coefficient's bit length
is a factor of $5^{39}$. It remains an enormous constant. This is
a reduction in coefficient height, not a practical claim that the
expanded polynomial can be stored.

\section{The ZFC equivalence and its verification status}
\label{sec:logical}

\subsection{Logical calibration}

At the pinned snapshot, the Lean source \src{KEndpoint215.lean}
states the exact endpoint
\begin{lstlisting}
theorem blank_halts_iff_zfc_refutation :
  BlankHalts machine215 <-> BProv ZFCax_s [] Form.fBot
\end{lstlisting}
and obtains its nonhalting counterpart by taking negations
\cite{endpoint}. Here \src{BProv} names the repository's finitary
proof relation; this endpoint is not a bound on the length of the
refutation. The formal source calibrates the raw table used here.

All the newly used arguments are ordinary finite mathematics available
in ZFC: arithmetic circuits, finite tableaux, Pell recurrences,
elementary divisibility and bounds, finite Boolean subset constructions,
and the classical three-square theorem for the integer-domain conversion.
Their assertions about witness existence, integer signs, and the fixed
machine can be arithmetized inside ZFC. The machine simulation and proof
enumeration establish the ordinary effective equivalence between
blank halting and existence of a finite ZFC refutation. Changing between
the repository's axiom enumeration and a conventional effective
presentation of ZFC does not change that consistency statement: the
finitary proof translations are provable in ZFC.

The proof is therefore the following chain, in the stated domain:
\begin{align*}
 \exists\mathbf x\ P(\mathbf x)=0
 &\ \longleftrightarrow\
 \text{the four-site native masks encode a valid blank tableau}\\
 &\ \longleftrightarrow\ \mathrm{BlankHalts}(\mathrm{K215})\\
 &\ \longleftrightarrow\ \mathrm{Ref}_{\mathrm{ZFC}}.
\end{align*}
Taking negations yields Theorem~\ref{thm:main}. This does not infer
ZFC provability from the soundness of a proof assistant: the
mathematical reductions themselves are available in ZFC. The inherited
formal endpoint documents the machine calibration; it is not a
newly compiled formalization of the changed Diophantine reductions.

There is no claim here to prove $\Con(\mathrm{ZFC})$. Subject to
the usual consistency assumption and the standard hypotheses of
the second incompleteness theorem, ZFC cannot prove the asserted
absence of solutions. The equivalence is a theorem about which
arithmetic assertion expresses that consistency statement.
No accepting run of the literal K215 machine, and no solution of
the final equation, was found or used.

\subsection{What was checked, and what remains mathematical}

The archive makes the following distinctions explicit.
\begin{enumerate}
\item \textbf{Inherited formal endpoint.}
The exact table and the statement of the Lean endpoint were inspected
at the pinned commit. The whole Lean/Rocq project was not rebuilt
in this work, and no claim of a fresh kernel check is made.
\item \textbf{New mathematical arguments.}
Weighted root restoration, the signed quotient extension, the
ordinary-strong combination, the smaller quadratic lift, the
ten-coordinate positivity reduction, and both smaller neighborhood
compilers are proved in this report and supporting notes.
The delicate denominator and signed-domain arguments received
independent review.
\item \textbf{Exact finite computation.}
Standard-library Python scripts check arithmetic-circuit dependency
graphs, full register liveness, operation counts, exact sparse
identities and residual degrees, the literal transition-table hash,
the local alphabet counts by exhaustive center categories, and
sparse native layout identities.
\item \textbf{Supplementary diagnostics.}
Random signed assignments and modular line restrictions check emitted
circuits against independently written formulas. Toy halting tableaux,
mutated periodic planes, and isolated root-restoration examples check
implementations. They are not substitutes for the general proofs or
witnesses for the final consistency equation.
\end{enumerate}
The ordinary strong rank lemma and the native mask compiler are
substantial inherited dependencies. Their precise hypotheses and
uses are recorded here and in the included upstream audit notes.
The new results should be reviewed as a mathematical contribution
with reproducible checks, not advertised as already formalized.

\section{Archive layout and reproduction}
\label{sec:reproduction}

The root \src{README.md} gives the recommended integration entry points.
The principal artifacts are:
\begin{itemize}
\item \src{con_zfc_reductions.tex} and its compiled PDF: this report.
\item \src{arithmetic/polynomial13_84.json} and
\src{arithmetic/polynomial13_lowweight85.json}: the complete live
circuits for the two thirteen-variable equations.
\item \src{arithmetic/system13_78.json},
\src{arithmetic/amalgam14_system67.json}, and associated product/SOS
packets: the lower-variable alternatives.
\item \src{low_degree/ordinary15_system61.json} and
\src{low_degree/retained16_system60.json}: full comparison circuits
and their single-polynomial forms.
\item \src{low_degree/quartic32.json} and
\src{low_degree/integer_quartic62.json}: all expanded quadratic
residuals defining the positive and unrestricted quartics.
\item \src{compiler/tshape_constants.json} and its generator:
the complete four-site coefficient recipe and evaluated scalar geometry.
The cross version provides a simpler intermediate compiler comparison.
\item \src{proofs/}: derivation notes and independent reviews.
\src{upstream/}: pinned dependencies, machine endpoint, source
comparison, and provenance. \src{SOURCE_PINS.json} and the hash
manifest identify the distributed source files.
\end{itemize}

From the extracted archive directory run
\begin{lstlisting}
python3 verify_arithmetic.py
python3 verify_compiler.py
latexmk -pdf -interaction=nonstopmode -halt-on-error con_zfc_reductions.tex
\end{lstlisting}
The mathematical checks require only Python's standard library and
local files. Typesetting requires a conventional TeX installation
with the packages named in the preamble. No network access, external
algebra service, giant coefficient materialization, or search for
a ZFC refutation is needed.

Operation counts refer to the archived circuits with their four
coefficient ports. These ports are bound by the new recipe in
\eqref{eq:ports}; they are not universally quantified parameters.
Some inherited circuit source packets retain their original
coefficient labels for provenance. Renaming their four ports
does not alter the formulas or evaluation cost. The manifest
distinguishes inherited inputs from newly emitted constructions.

\section{Conclusions and sensible next targets}

For incorporation, the $13$-variable, $84$-operation product is the
recommended compact equation, with the degree-$515$ alternative
available at one additional operation. If low degree is primary,
the $32$-positive-variable quartic and $62$-integer-variable quartic
are the useful outputs. If evaluation cost dominates, the old
$67$-operation product remains competitive; the new $61$-operation
system improves its degree without lowering its cost.

The gains arise at different layers and compose cleanly:
ordinary strong scaling changes an existential Pell multiplier;
weighted root amalgamation removes witnesses while guarding the
zero-denominator branch; register elimination lowers the quartic's
variable count; translating one local test removes a neighborhood
coordinate and shrinks the native compiler. None requires changing
the underlying proof-search machine.

Natural further work is to formalize the weighted-root lemma and
signed-domain extension, then the three-overlap compiler, before
attempting new optimization claims. A search over still smaller
tile presentations or a different exact coefficient layout could
improve coefficient height again. This report establishes concrete
improvements over the reviewed packages, not a theorem that any
one displayed measure is globally minimal.

\begin{thebibliography}{99}
\raggedright
\bibitem{catalogue}
V. Reshetnikov and contributors, \emph{Con(ZFC) in Diophantine form:
comparison of the eight report packages}, ProveIt repository,
snapshot \pin.
\href{https://github.com/VladimirReshetnikov/ProveIt/tree/3b201cdea32d2c816f5f3d20734051f25fd3ef14/SetTheory/Cardinals/docs/reports/hilbert-tenth-problem/con-zfc-diophantine-forms}{Pinned comparison and packages 01--08}.

\bibitem{blank}
ProveIt, package 06, \emph{The 67-operation blank-input polynomial},
including \src{report.tex}, \src{proofs/fixed_blank_67_proof.md},
\src{proofs/pell_lemma_audit.md}, and \src{code/blank_helical.py}.
\href{https://github.com/VladimirReshetnikov/ProveIt/tree/3b201cdea32d2c816f5f3d20734051f25fd3ef14/SetTheory/Cardinals/docs/reports/hilbert-tenth-problem/con-zfc-diophantine-forms/06-blank67}{Pinned package 06}.

\bibitem{ordinary}
ProveIt, package 02, \src{FIXED_INPUT75_ORDINARY_PROOF.md} and
\src{PELL_INTERFACE_LEMMAS.md}.
\href{https://github.com/VladimirReshetnikov/ProveIt/tree/3b201cdea32d2c816f5f3d20734051f25fd3ef14/SetTheory/Cardinals/docs/reports/hilbert-tenth-problem/con-zfc-diophantine-forms/02-system63}{Pinned ordinary-strong package}.

\bibitem{relaxed}
ProveIt, \src{Papers/1980/PELL_RELAXED_AUXILIARY_PROOF.md},
ordinary strong rank lemma and Pell divisibility.
\href{https://github.com/VladimirReshetnikov/ProveIt/blob/3b201cdea32d2c816f5f3d20734051f25fd3ef14/Computability/HilbertTenthProblem/Papers/1980/PELL_RELAXED_AUXILIARY_PROOF.md}{Pinned proof note}.

\bibitem{signed}
ProveIt, \src{complete84_signed_quotient_absorption.md},
signed auxiliary-root step-down.
\href{https://github.com/VladimirReshetnikov/ProveIt/blob/3b201cdea32d2c816f5f3d20734051f25fd3ef14/Computability/HilbertTenthProblem/Papers/research-wip/native-stream-queue/complete84_signed_quotient_absorption.md}{Pinned signed-quotient note}.

\bibitem{scout}
ProveIt, \src{complete84_signed_quotient_soundness_scout.md},
sign-independent forward power, extraction, and population arguments.
\href{https://github.com/VladimirReshetnikov/ProveIt/blob/3b201cdea32d2c816f5f3d20734051f25fd3ef14/Computability/HilbertTenthProblem/Papers/research-wip/native-stream-queue/complete84_signed_quotient_soundness_scout.md}{Pinned soundness note}.

\bibitem{half}
ProveIt, \src{pell_kernel_half_binomial42.md},
strong half-binomial Pell kernel.
\href{https://github.com/VladimirReshetnikov/ProveIt/blob/3b201cdea32d2c816f5f3d20734051f25fd3ef14/Computability/HilbertTenthProblem/Papers/research-wip/native-stream-queue/pell_kernel_half_binomial42.md}{Pinned kernel note}.

\bibitem{helical}
ProveIt, package 06, \src{proofs/blank_helical_proof.md},
finite-strip soundness and padded helical completeness.
\href{https://github.com/VladimirReshetnikov/ProveIt/tree/3b201cdea32d2c816f5f3d20734051f25fd3ef14/SetTheory/Cardinals/docs/reports/hilbert-tenth-problem/con-zfc-diophantine-forms/06-blank67/proofs}{Pinned compiler proofs}.

\bibitem{endpoint}
ProveIt, \src{KEndpoint215.lean} and the K215 verification package,
\src{blank_halts_iff_zfc_refutation} and
\src{nonhalting_iff_zfc_consistency}.
\href{https://github.com/VladimirReshetnikov/ProveIt/tree/3b201cdea32d2c816f5f3d20734051f25fd3ef14/Computability/BusyBeaver/Research/ZFC/k225-formal/verification/k215}{Pinned K215 verification}.
The local endpoint source is included in \src{upstream/logic}.

\bibitem{nine}
ProveIt, \src{Diophantine/Paper1982/NineMain.lean},
\src{Jones1982.nine_unknowns}, and its formalized dependencies.
\href{https://github.com/VladimirReshetnikov/ProveIt/blob/3b201cdea32d2c816f5f3d20734051f25fd3ef14/Computability/HilbertTenthProblem/Lean/Diophantine/Paper1982/NineMain.lean}{Pinned nine-unknowns existence theorem}.
\end{thebibliography}
\end{document}



